\maketitle

\thispagestyle{empty}
\vspace*{\fill}
\noindent\small
Copyright \textcopyright{} 2026 Bijan Fallah and the Chronos-ESM contributors.\\
This manual accompanies the Chronos-ESM source code (Apache License 2.0) and is intended
for research and teaching. It is a living document, generated alongside the code so that
every governing equation is traceable to its implementation.\\[1em]
Source: \url{https://github.com/bijanf/ESMB}
\vspace*{\fill}
\clearpage

\chapter*{Preface}
\addcontentsline{toc}{chapter}{Preface}

\model{} is a fully \emph{differentiable}, GPU- and CPU-capable coupled Earth System Model
written in \jax{}. It couples a multi-level spectral primitive-equation atmosphere (built on
Google's \dino{} dynamical core), a $z$-level primitive-equation ocean, a thermodynamic sea
ice model, and a land-surface scheme, at a teaching-friendly T31 ($\approx 3.75^\circ$)
resolution that runs in minutes-to-hours on a single GPU.

This book has three aims. First, to be \textbf{equation-complete and honest}: every
prognostic equation, closure, and numerical scheme is derived from, and cross-referenced to,
the actual source code --- including its approximations and known biases. Climate-model
documentation too often hides the gap between the textbook equations and what the code really
integrates; here the two are presented together, with the code location given beside each
equation. Second, to be \textbf{pedagogical}: the model is small and fast enough that a
student can run a coupled climate experiment, change an equation, and see the result the same
afternoon. Each chapter ends with exercises that use the model as a laboratory. Third, to
showcase \textbf{differentiable Earth-system modeling}: because the entire model is written
in \jax{}, one can take the gradient of any output (global temperature, AMOC strength, a
tipping threshold) with respect to any input or parameter --- enabling calibration,
sensitivity analysis, and data assimilation that are impractical with conventional models.

The book is organized in five parts: \emph{Foundations} (the continuous and discrete
mathematics, and the differentiable-programming paradigm); \emph{The Component Models}
(ocean, atmosphere, ice, land, and the coupler, each derived from the code); \emph{Forcing,
Sensitivity, and Data Assimilation}; \emph{Canonical Climate Test Cases} (the standard
phenomena every climate model must reproduce --- mean climatology, the AMOC and its tipping,
mid-latitude jets and storm tracks, the ITCZ and monsoons); and \emph{Community, Standards,
and Teaching}, including a roadmap toward CMIP7 participation and a guide to using \model{} in
the classroom.

\bigskip
\noindent\textit{A note on status.} \model{} is research software under active development.
Where a component is a deliberate simplification or a work in progress, this book says so
explicitly --- the honest documentation of an evolving model is part of its scientific value.
